% Informe del TP de Implementación: Esteganografía (Criptografía y Seguridad, 72.04, ITBA).
% Se compila con `make report` (o `bash informe/build.sh`); no usar latexmk (no está instalado).
% Las cifras medidas de las Cuestiones II y VII salen de generado/ (regenerado por medir_q2.py).
\documentclass[a4paper,11pt]{article}

\usepackage[utf8]{inputenc}
\usepackage[T1]{fontenc}
\usepackage{lmodern}
% El paquete de hyphenation español no está instalado en todos los entornos: se importa el idioma con
% \babelprovide (la opción `spanish` de babel falla aquí con "Unknown option").
\usepackage{babel}
\babelprovide[import,main]{spanish}
\usepackage[a4paper,margin=2.5cm]{geometry}
\usepackage{microtype}
\usepackage{booktabs}
\usepackage{tabularx}
\usepackage{array}
\usepackage{multirow}
\usepackage{graphicx}
\usepackage{subcaption}
\usepackage{float}
\usepackage{xcolor}
\usepackage{amsmath}
\usepackage{listings}
\usepackage[hidelinks,unicode=true]{hyperref}
\input{glyphtounicode}
\pdfgentounicode=1

\setcounter{secnumdepth}{0}
\setlength{\emergencystretch}{3em}
\setlength{\parindent}{0pt}
\setlength{\parskip}{0.6em}
\lstset{basicstyle=\ttfamily\small,breaklines=true,columns=fullflexible,keepspaces=true,
  frame=single,framerule=0.3pt,xleftmargin=0.5em,xrightmargin=0.5em,aboveskip=0.8em,belowskip=0.8em}

% ---------------------------------------------------------------- datos medidos
% \definirdato{clave}{valor} guarda un valor calculado; \dato{clave} lo imprime y detiene la
% compilación si la clave no existe (una medición faltante nunca pasa inadvertida).
\makeatletter
\newcommand{\definirdato}[2]{\expandafter\gdef\csname dato@#1\endcsname{#2}}
\newcommand{\dato}[1]{%
  \ifcsname dato@#1\endcsname
    \csname dato@#1\endcsname
  \else
    \PackageError{informe}{Dato no definido: #1}{Falta la clave en generado/datos.tex; ejecutar make report.}%
  \fi}
\makeatother
% Generado por informe/scripts/medir_q2.py: no editar a mano.
% Claves definidas (\dato{clave}); tex_int agrupa miles con \, y los decimales usan {,}:
%   portador-ancho
%   portador-alto
%   portador-bytes
%   cap-LSB1
%   cap-LSB4
%   cap-LSBI
%   png-bytes
%   png-LSB1-ocultos
%   png-LSB1-uso
%   png-LSB1-usados
%   png-LSB1-cambiados
%   png-LSB1-bits
%   png-LSB1-tasa
%   png-LSB1-pctportador
%   png-LSB1-mse
%   png-LSB1-psnr
%   png-LSB1-rojos
%   png-LSB1-ultimo
%   png-LSB1-cambiomax
%   png-LSB4-ocultos
%   png-LSB4-uso
%   png-LSB4-usados
%   png-LSB4-cambiados
%   png-LSB4-bits
%   png-LSB4-tasa
%   png-LSB4-pctportador
%   png-LSB4-mse
%   png-LSB4-psnr
%   png-LSB4-rojos
%   png-LSB4-ultimo
%   png-LSB4-cambiomax
%   png-LSBI-ocultos
%   png-LSBI-uso
%   png-LSBI-usados
%   png-LSBI-cambiados
%   png-LSBI-bits
%   png-LSBI-tasa
%   png-LSBI-pctportador
%   png-LSBI-mse
%   png-LSBI-psnr
%   png-LSBI-rojos
%   png-LSBI-ultimo
%   png-LSBI-cambiomax
%   paper-bytes
%   paper-LSB1-ocultos
%   paper-LSB1-uso
%   paper-LSB1-usados
%   paper-LSB1-cambiados
%   paper-LSB1-bits
%   paper-LSB1-tasa
%   paper-LSB1-pctportador
%   paper-LSB1-mse
%   paper-LSB1-psnr
%   paper-LSB1-rojos
%   paper-LSB1-ultimo
%   paper-LSB1-cambiomax
%   paper-LSB4-ocultos
%   paper-LSB4-uso
%   paper-LSB4-usados
%   paper-LSB4-cambiados
%   paper-LSB4-bits
%   paper-LSB4-tasa
%   paper-LSB4-pctportador
%   paper-LSB4-mse
%   paper-LSB4-psnr
%   paper-LSB4-rojos
%   paper-LSB4-ultimo
%   paper-LSB4-cambiomax
%   paper-LSBI-ocultos
%   paper-LSBI-uso
%   paper-LSBI-usados
%   paper-LSBI-cambiados
%   paper-LSBI-bits
%   paper-LSBI-tasa
%   paper-LSBI-pctportador
%   paper-LSBI-mse
%   paper-LSBI-psnr
%   paper-LSBI-rojos
%   paper-LSBI-ultimo
%   paper-LSBI-cambiomax
%   azar-bytes
%   azar-LSB1-ocultos
%   azar-LSB1-uso
%   azar-LSB1-usados
%   azar-LSB1-cambiados
%   azar-LSB1-bits
%   azar-LSB1-tasa
%   azar-LSB1-pctportador
%   azar-LSB1-mse
%   azar-LSB1-psnr
%   azar-LSB1-rojos
%   azar-LSB1-ultimo
%   azar-LSB1-cambiomax
%   azar-LSB4-ocultos
%   azar-LSB4-uso
%   azar-LSB4-usados
%   azar-LSB4-cambiados
%   azar-LSB4-bits
%   azar-LSB4-tasa
%   azar-LSB4-pctportador
%   azar-LSB4-mse
%   azar-LSB4-psnr
%   azar-LSB4-rojos
%   azar-LSB4-ultimo
%   azar-LSB4-cambiomax
%   azar-LSBI-ocultos
%   azar-LSBI-uso
%   azar-LSBI-usados
%   azar-LSBI-cambiados
%   azar-LSBI-bits
%   azar-LSBI-tasa
%   azar-LSBI-pctportador
%   azar-LSBI-mse
%   azar-LSBI-psnr
%   azar-LSBI-rojos
%   azar-LSBI-ultimo
%   azar-LSBI-cambiomax
%   est-portador-B-lsb
%   est-portador-B-chi
%   est-portador-G-lsb
%   est-portador-G-chi
%   est-portador-R-lsb
%   est-portador-R-chi
%   est-png-LSB1-B-lsb
%   est-png-LSB1-B-chi
%   est-png-LSB1-G-lsb
%   est-png-LSB1-G-chi
%   est-png-LSB1-R-lsb
%   est-png-LSB1-R-chi
%   est-png-LSB4-B-lsb
%   est-png-LSB4-B-chi
%   est-png-LSB4-G-lsb
%   est-png-LSB4-G-chi
%   est-png-LSB4-R-lsb
%   est-png-LSB4-R-chi
%   est-png-LSBI-B-lsb
%   est-png-LSBI-B-chi
%   est-png-LSBI-G-lsb
%   est-png-LSBI-G-chi
%   est-png-LSBI-R-lsb
%   est-png-LSBI-R-chi
%   inv-png-n
%   inv-png-sin
%   inv-png-con
%   inv-png-reduccion
%   inv-png-flags
%   inv-png-banderas
%   inv-paper-n
%   inv-paper-sin
%   inv-paper-con
%   inv-paper-reduccion
%   inv-paper-flags
%   inv-paper-banderas
%   inv-azar-n
%   inv-azar-sin
%   inv-azar-con
%   inv-azar-reduccion
%   inv-azar-flags
%   inv-azar-banderas
%   inv-aes256ofb-n
%   inv-aes256ofb-sin
%   inv-aes256ofb-con
%   inv-aes256ofb-reduccion
%   inv-aes256ofb-flags
%   inv-aes256ofb-banderas
%   inv-descfb-n
%   inv-descfb-sin
%   inv-descfb-con
%   inv-descfb-reduccion
%   inv-descfb-flags
%   inv-descfb-banderas
%   inv-azar-teorico
%   inv-azar-teoricopct
\definirdato{portador-ancho}{640}
\definirdato{portador-alto}{480}
\definirdato{portador-bytes}{921\,600}
\definirdato{cap-LSB1}{115\,200}
\definirdato{cap-LSB4}{460\,800}
\definirdato{cap-LSBI}{76\,799}
\definirdato{png-bytes}{44\,886}
\definirdato{png-LSB1-ocultos}{44\,895}
\definirdato{png-LSB1-uso}{38{,}97}
\definirdato{png-LSB1-usados}{359\,160}
\definirdato{png-LSB1-cambiados}{190\,182}
\definirdato{png-LSB1-bits}{190\,182}
\definirdato{png-LSB1-tasa}{52{,}95}
\definirdato{png-LSB1-pctportador}{20{,}64}
\definirdato{png-LSB1-mse}{0{,}2064}
\definirdato{png-LSB1-psnr}{54{,}98}
\definirdato{png-LSB1-rojos}{63\,341}
\definirdato{png-LSB1-ultimo}{359\,213}
\definirdato{png-LSB1-cambiomax}{1}
\definirdato{png-LSB4-ocultos}{44\,895}
\definirdato{png-LSB4-uso}{9{,}74}
\definirdato{png-LSB4-usados}{89\,790}
\definirdato{png-LSB4-cambiados}{85\,370}
\definirdato{png-LSB4-bits}{205\,965}
\definirdato{png-LSB4-tasa}{95{,}08}
\definirdato{png-LSB4-pctportador}{9{,}26}
\definirdato{png-LSB4-mse}{8{,}7631}
\definirdato{png-LSB4-psnr}{38{,}70}
\definirdato{png-LSB4-rojos}{28\,443}
\definirdato{png-LSB4-ultimo}{89\,843}
\definirdato{png-LSB4-cambiomax}{15}
\definirdato{png-LSBI-ocultos}{44\,895}
\definirdato{png-LSBI-uso}{58{,}46}
\definirdato{png-LSBI-usados}{359\,164}
\definirdato{png-LSBI-cambiados}{170\,852}
\definirdato{png-LSBI-bits}{170\,852}
\definirdato{png-LSBI-tasa}{47{,}57}
\definirdato{png-LSBI-pctportador}{18{,}54}
\definirdato{png-LSBI-mse}{0{,}1854}
\definirdato{png-LSBI-psnr}{55{,}45}
\definirdato{png-LSBI-rojos}{0}
\definirdato{png-LSBI-ultimo}{538\,794}
\definirdato{png-LSBI-cambiomax}{1}
\definirdato{paper-bytes}{226}
\definirdato{paper-LSB1-ocultos}{235}
\definirdato{paper-LSB1-uso}{0{,}20}
\definirdato{paper-LSB1-usados}{1\,880}
\definirdato{paper-LSB1-cambiados}{1\,051}
\definirdato{paper-LSB1-bits}{1\,051}
\definirdato{paper-LSB1-tasa}{55{,}90}
\definirdato{paper-LSB1-pctportador}{0{,}11}
\definirdato{paper-LSB1-mse}{0{,}0011}
\definirdato{paper-LSB1-psnr}{77{,}56}
\definirdato{paper-LSB1-rojos}{354}
\definirdato{paper-LSB1-ultimo}{1\,933}
\definirdato{paper-LSB1-cambiomax}{1}
\definirdato{paper-LSB4-ocultos}{235}
\definirdato{paper-LSB4-uso}{0{,}05}
\definirdato{paper-LSB4-usados}{470}
\definirdato{paper-LSB4-cambiados}{454}
\definirdato{paper-LSB4-bits}{1\,041}
\definirdato{paper-LSB4-tasa}{96{,}60}
\definirdato{paper-LSB4-pctportador}{0{,}05}
\definirdato{paper-LSB4-mse}{0{,}0527}
\definirdato{paper-LSB4-psnr}{60{,}91}
\definirdato{paper-LSB4-rojos}{150}
\definirdato{paper-LSB4-ultimo}{523}
\definirdato{paper-LSB4-cambiomax}{15}
\definirdato{paper-LSBI-ocultos}{235}
\definirdato{paper-LSBI-uso}{0{,}31}
\definirdato{paper-LSBI-usados}{1\,884}
\definirdato{paper-LSBI-cambiados}{835}
\definirdato{paper-LSBI-bits}{835}
\definirdato{paper-LSBI-tasa}{44{,}32}
\definirdato{paper-LSBI-pctportador}{0{,}09}
\definirdato{paper-LSBI-mse}{0{,}0009}
\definirdato{paper-LSBI-psnr}{78{,}56}
\definirdato{paper-LSBI-rojos}{1}
\definirdato{paper-LSBI-ultimo}{2\,862}
\definirdato{paper-LSBI-cambiomax}{1}
\definirdato{azar-bytes}{40\,000}
\definirdato{azar-LSB1-ocultos}{40\,009}
\definirdato{azar-LSB1-uso}{34{,}73}
\definirdato{azar-LSB1-usados}{320\,072}
\definirdato{azar-LSB1-cambiados}{159\,834}
\definirdato{azar-LSB1-bits}{159\,834}
\definirdato{azar-LSB1-tasa}{49{,}94}
\definirdato{azar-LSB1-pctportador}{17{,}34}
\definirdato{azar-LSB1-mse}{0{,}1734}
\definirdato{azar-LSB1-psnr}{55{,}74}
\definirdato{azar-LSB1-rojos}{53\,365}
\definirdato{azar-LSB1-ultimo}{320\,125}
\definirdato{azar-LSB1-cambiomax}{1}
\definirdato{azar-LSB4-ocultos}{40\,009}
\definirdato{azar-LSB4-uso}{8{,}68}
\definirdato{azar-LSB4-usados}{80\,018}
\definirdato{azar-LSB4-cambiados}{74\,983}
\definirdato{azar-LSB4-bits}{160\,169}
\definirdato{azar-LSB4-tasa}{93{,}71}
\definirdato{azar-LSB4-pctportador}{8{,}14}
\definirdato{azar-LSB4-mse}{5{,}7478}
\definirdato{azar-LSB4-psnr}{40{,}54}
\definirdato{azar-LSB4-rojos}{24\,977}
\definirdato{azar-LSB4-ultimo}{80\,071}
\definirdato{azar-LSB4-cambiomax}{15}
\definirdato{azar-LSBI-ocultos}{40\,009}
\definirdato{azar-LSBI-uso}{52{,}10}
\definirdato{azar-LSBI-usados}{320\,076}
\definirdato{azar-LSBI-cambiados}{159\,773}
\definirdato{azar-LSBI-bits}{159\,773}
\definirdato{azar-LSBI-tasa}{49{,}92}
\definirdato{azar-LSBI-pctportador}{17{,}34}
\definirdato{azar-LSBI-mse}{0{,}1734}
\definirdato{azar-LSBI-psnr}{55{,}74}
\definirdato{azar-LSBI-rojos}{1}
\definirdato{azar-LSBI-ultimo}{480\,163}
\definirdato{azar-LSBI-cambiomax}{1}
\definirdato{est-portador-B-lsb}{0{,}5475}
\definirdato{est-portador-B-chi}{8475{,}17}
\definirdato{est-portador-G-lsb}{0{,}5442}
\definirdato{est-portador-G-chi}{7592{,}63}
\definirdato{est-portador-R-lsb}{0{,}5488}
\definirdato{est-portador-R-chi}{8606{,}07}
\definirdato{est-png-LSB1-B-lsb}{0{,}4582}
\definirdato{est-png-LSB1-B-chi}{2125{,}38}
\definirdato{est-png-LSB1-G-lsb}{0{,}4571}
\definirdato{est-png-LSB1-G-chi}{2144{,}94}
\definirdato{est-png-LSB1-R-lsb}{0{,}4577}
\definirdato{est-png-LSB1-R-chi}{2028{,}76}
\definirdato{est-png-LSB4-B-lsb}{0{,}5137}
\definirdato{est-png-LSB4-B-chi}{5232{,}30}
\definirdato{est-png-LSB4-G-lsb}{0{,}5108}
\definirdato{est-png-LSB4-G-chi}{4372{,}81}
\definirdato{est-png-LSB4-R-lsb}{0{,}5147}
\definirdato{est-png-LSB4-R-chi}{5226{,}22}
\definirdato{est-png-LSBI-B-lsb}{0{,}5626}
\definirdato{est-png-LSBI-B-chi}{3520{,}98}
\definirdato{est-png-LSBI-G-lsb}{0{,}5608}
\definirdato{est-png-LSBI-G-chi}{3345{,}95}
\definirdato{est-png-LSBI-R-lsb}{0{,}5488}
\definirdato{est-png-LSBI-R-chi}{8606{,}07}
\definirdato{inv-png-n}{359\,160}
\definirdato{inv-png-sin}{188\,308}
\definirdato{inv-png-con}{170\,852}
\definirdato{inv-png-reduccion}{9{,}27}
\definirdato{inv-png-flags}{1111}
\definirdato{inv-png-banderas}{0}
\definirdato{inv-paper-n}{1\,880}
\definirdato{inv-paper-sin}{1\,044}
\definirdato{inv-paper-con}{832}
\definirdato{inv-paper-reduccion}{20{,}31}
\definirdato{inv-paper-flags}{0001}
\definirdato{inv-paper-banderas}{3}
\definirdato{inv-azar-n}{320\,072}
\definirdato{inv-azar-sin}{160\,032}
\definirdato{inv-azar-con}{159\,771}
\definirdato{inv-azar-reduccion}{0{,}16}
\definirdato{inv-azar-flags}{0101}
\definirdato{inv-azar-banderas}{2}
\definirdato{inv-aes256ofb-n}{359\,192}
\definirdato{inv-aes256ofb-sin}{179\,491}
\definirdato{inv-aes256ofb-con}{179\,329}
\definirdato{inv-aes256ofb-reduccion}{0{,}09}
\definirdato{inv-aes256ofb-flags}{0001}
\definirdato{inv-aes256ofb-banderas}{3}
\definirdato{inv-descfb-n}{359\,192}
\definirdato{inv-descfb-sin}{179\,425}
\definirdato{inv-descfb-con}{179\,195}
\definirdato{inv-descfb-reduccion}{0{,}13}
\definirdato{inv-descfb-flags}{1000}
\definirdato{inv-descfb-banderas}{3}
\definirdato{inv-azar-teorico}{439}
\definirdato{inv-azar-teoricopct}{0{,}27}


% ---------------------------------------------------------------- estructura
\newenvironment{enunciado}{\begin{quote}\itshape\noindent\textbf{Enunciado.}\ }{\end{quote}}
\newcommand{\pendienteEstegoanalisis}{%
  \begin{center}\fbox{\parbox{0.9\linewidth}{\itshape Pendiente: esta cuestión requiere los archivos de estegoanálisis de la cátedra, que todavía no fueron entregados.}}\end{center}}
\newcommand{\respuesta}[1]{\IfFileExists{secciones/#1.tex}{\input{secciones/#1}}{}}

\title{Trabajo Práctico de Implementación: Esteganografía}
\author{Lucas Di Candia}
\date{\today}
\hypersetup{pdftitle={Trabajo Práctico de Implementación: Esteganografía},pdfauthor={Lucas Di Candia}}

\begin{document}

\begin{titlepage}
  \centering
  \vspace*{4cm}
  {\Huge\bfseries Trabajo Práctico de Implementación:\\Esteganografía\par}
  \vspace{1.5cm}
  {\Large Criptografía y Seguridad (72.04) --- Instituto Tecnológico de Buenos Aires\par}
  \vspace{2cm}
  {\Large Lucas Di Candia\par}
  \vspace{1cm}
  {\large\today\par}
\end{titlepage}

\tableofcontents
\clearpage

% ---------------------------------------------------------------- introducción
\section{Introducción}\label{sec:introduccion}

\texttt{stegobmp} es un programa de línea de comandos escrito en C, con OpenSSL, que oculta cualquier
archivo dentro de una imagen BMP de 24 bits sin compresión (versión 3) y lo recupera después. Implementa
los tres métodos que pide la cátedra~\cite{catedra2026}: LSB1, LSB4 y LSBI, este último la variante con
inversión de bits de Majeed y Sulaiman~\cite{majeed2015}. Opcionalmente, antes de ocultar, cifra el
mensaje con AES (128, 192 o 256 bits) o con 3DES, en los modos ECB, CFB, OFB o CBC, derivando la clave
y el vector de inicialización con PBKDF2 a partir de una contraseña.

La esteganografía oculta la existencia misma del mensaje, a diferencia de la criptografía, que oculta su contenido~\cite{johnson1998,gomezcardenas}; por eso el programa permite combinar ambas.

\textbf{Organización.} El informe responde en orden las cuestiones de la sección 7 del enunciado, una sección por cuestión, y cada respuesta empieza
transcribiendo la pregunta. Las Cuestiones I y II analizan el documento de Majeed y Sulaiman y comparan los tres métodos; las Cuestiones VII y VIII explican y discuten la mejora que propone LSBI; la IX cuenta las
dificultades de la implementación y la X propone mejoras. Las Cuestiones III a VI dependen de los archivos de estegoanálisis de la cátedra, que todavía no fueron entregados, y
quedan marcadas como pendientes.

\textbf{Convenciones.} Las páginas citadas del documento de Majeed y Sulaiman son las de la revista (342 a 348). Todas las cifras medidas que se
citan en las Cuestiones II y VII (capacidades, bytes modificados, MSE, PSNR, estadísticos $\chi^2$, conteos de inversión) se
regeneran con \texttt{make report} a partir del propio programa y de las herramientas de análisis del
repositorio, y se contrastaron con los archivos de \texttt{Ejemplo/} de la cátedra; ninguna se escribió a mano. Las rutas como \texttt{docs/LSBI-NOTES.md} remiten a las notas del proyecto.
Para regenerar este informe basta ejecutar \texttt{make report} en la raíz del repositorio.

% ---------------------------------------------------------------- cuestiones
\section{Cuestión I: Análisis del documento de Majeed y Sulaiman}\label{sec:cuestion-i}
\begin{enunciado}
Discutir los siguientes aspectos relativos al documento. a) Organización formal del documento. b) La descripción del algoritmo. c) La notación utilizada, ¿es clara? ¿hay algún error o contradicción?
\end{enunciado}
\respuesta{q01-paper}

\section{Cuestión II: Comparación de LSB1, LSB4 y LSBI}\label{sec:cuestion-ii}
\begin{enunciado}
Esteganografiar un mismo archivo en un .bmp con cada uno de los tres algoritmos, y comparar los resultados obtenidos. Hacer un cuadro comparativo de los tres algoritmos estableciendo ventajas y desventajas.
\end{enunciado}
\respuesta{q02-comparacion}

\section{Cuestión III: Procedimiento de estegoanálisis}\label{sec:cuestion-iii}
\begin{enunciado}
Explicar detalladamente el procedimiento realizado para descubrir qué se había ocultado en cada archivo y de qué modo. Indicar qué se encontró en cada archivo.
\end{enunciado}
\pendienteEstegoanalisis

\section{Cuestión IV: Mensajes ocultos dentro de mensajes ocultos}\label{sec:cuestion-iv}
\begin{enunciado}
Algunos mensajes ocultos tenían, a su vez, otros mensajes ocultos. Indica cuál era ese mensaje y cómo se había ocultado.
\end{enunciado}
\pendienteEstegoanalisis

\section{Cuestión V: El fragmento de película}\label{sec:cuestion-v}
\begin{enunciado}
Uno de los archivos ocultos era una porción de un video de una película, donde se ve ejemplificado una manera de ocultar información ¿qué se ocultaba y sobre qué portador?
\end{enunciado}
\pendienteEstegoanalisis

\section{Cuestión VI: El método de ocultamiento desconocido}\label{sec:cuestion-vi}
\begin{enunciado}
¿De qué se trató el método de esteganografiado que no era LSB1 ni LSB4 ni LSBI? ¿Es un método eficaz? ¿Por qué?
\end{enunciado}
\pendienteEstegoanalisis

\section{Cuestión VII: Por qué LSBI mejora a LSB}\label{sec:cuestion-vii}
\begin{enunciado}
¿por qué la propuesta del documento de Majeed y Sulaiman es realmente una mejora respecto de LSB común?
\end{enunciado}
\respuesta{q07-mejora-lsbi}

\section{Cuestión VIII: Otras formas de guardar los patrones invertidos}\label{sec:cuestion-viii}
\begin{enunciado}
En la implementación se optó por guardar los patrones invertidos antes del mensaje ¿de qué otra manera o en qué otro lugar podría guardarse el registro de los patrones invertidos?
\end{enunciado}
\respuesta{q08-registro-patrones}

\section{Cuestión IX: Dificultades de implementación}\label{sec:cuestion-ix}
\begin{enunciado}
¿Qué dificultades encontraron en la implementación del algoritmo del paper?
\end{enunciado}
\respuesta{q09-dificultades}

\section{Cuestión X: Mejoras y extensiones futuras}\label{sec:cuestion-x}
\begin{enunciado}
¿Qué mejoras o futuras extensiones harías al programa stegobmp?
\end{enunciado}
\respuesta{q10-mejoras}

% ---------------------------------------------------------------- referencias
\begin{thebibliography}{99}

\bibitem{majeed2015} M. A. Majeed and R. Sulaiman, ``An improved LSB image steganography technique using bit-inverse in 24 bit colour image'', \emph{Journal of Theoretical and Applied Information Technology}, vol. 80, no. 2, pp. 342--348, 2015.

\bibitem{catedra2026} Cátedra de Criptografía y Seguridad, ITBA, ``Trabajo Práctico de Implementación: Esteganografía'', 2026.

\bibitem{johnson1998} N. F. Johnson and S. Jajodia, ``Exploring steganography: Seeing the unseen'', \emph{IEEE Computer}, vol. 31, no. 2, pp. 26--34, 1998.

\bibitem{westfeld2000} A. Westfeld and A. Pfitzmann, ``Attacks on steganographic systems'', \emph{Information Hiding 1999}, LNCS 1768, pp. 61--76, Springer, 2000.

\bibitem{fridrich2001} J. Fridrich, M. Goljan and R. Du, ``Reliable detection of LSB steganography in color and grayscale images'', \emph{Proc. ACM Workshop on Multimedia and Security}, pp. 27--30, 2001.

\bibitem{fridrich2009} J. Fridrich, \emph{Steganography in Digital Media: Principles, Algorithms, and Applications}, Cambridge University Press, 2009.

\bibitem{sharp2001} T. Sharp, ``An implementation of key-based digital signal steganography'', \emph{Information Hiding 2001}, LNCS 2137, pp. 13--26, Springer, 2001.

\bibitem{kerckhoffs1883} A. Kerckhoffs, ``La cryptographie militaire'', \emph{Journal des sciences militaires}, vol. IX, pp. 5--38, 1883.

\bibitem{rfc8018} K. Moriarty, B. Kaliski and A. Rusch, ``PKCS \#5: Password-Based Cryptography Specification Version 2.1'', RFC 8018, 2017.

\bibitem{nist80038a} M. Dworkin, ``Recommendation for Block Cipher Modes of Operation: Methods and Techniques'', NIST SP 800-38A, 2001.

\bibitem{nist80038d} M. Dworkin, ``Recommendation for Block Cipher Modes of Operation: Galois/Counter Mode (GCM) and GMAC'', NIST SP 800-38D, 2007.

\bibitem{nist800132} M. Sönmez Turan, E. Barker, W. Burr and L. Chen, ``Recommendation for Password-Based Key Derivation, Part 1: Storage Applications'', NIST SP 800-132, 2010.

\bibitem{rfc9106} A. Biryukov, D. Dinu, D. Khovratovich and S. Josefsson, ``Argon2 Memory-Hard Function for Password Hashing and Proof-of-Work Applications'', RFC 9106, 2021.

\bibitem{microsoftbmp} Microsoft, ``BITMAPFILEHEADER structure (wingdi.h)'', documentación de la API Win32, enlace indicado en el enunciado del TP.

\bibitem{gomezcardenas} R. Gómez Cárdenas, ``Esteganografía'', \url{http://www.cryptomex.org/SlidesCripto/Estegano.pdf}.

\bibitem{opensslevp} OpenSSL Project, ``EVP\_EncryptInit'' (página de manual), \url{https://www.openssl.org/docs/}.

\end{thebibliography}

\end{document}
